
\section{Representation and forecasting procedure}\label{sec:method}
\subsection{Prediction target and notation}
Let $(y_t)$ be a scalar observed time series. At origin $t$, the available history is $x_t=(y_{t-d+1},\ldots,y_t)^\top\in\R^d$, and the horizon-$h$ target is $y_{t+h}$, where $h$ belongs to a finite set $\mathcal H$ of positive integers. We fit a separate direct forecast for each horizon. Vectors are columns, a superscript $\top$ denotes transpose, $\|\cdot\|_2$ is the Euclidean norm, and $\|\cdot\|_1$ is the sum of absolute coordinates. The matrix $I_n$ and vector $\mathbf1_n$ are the identity and all-ones vector of dimension $n$. The operator $\concat$ stacks its vector arguments in a fixed order.

Independent episode sets A, B, C, and D have distinct roles. A fits input transformations, reference geometry, feature transformations, and readout coefficients. B selects hyperparameters. C supplies diagnostics only and makes no fitting or selection decision. D is the terminal evaluation set. Let $z\in\R^d$ denote the history after coordinatewise standardization using A. A history standard deviation at most $10^{-12}$ is replaced by one. All algebraic comparisons below condition on the fitted reference geometry.

\subsection{Local reference geometry and memberships}
A reference bank contains local models indexed by $j$, grouped into partitions indexed by $b=1,\ldots,B_0$. Let $\mathcal J_b$ be the model indices in partition $b$. Model $j$ contains $n_j$ standardized reference histories $q_{ji}\in\R^d$, stored as the rows $q_{ji}^\top$ of $Q_j\in\R^{n_j\times d}$. Its reference mean is $\mu_j$, and $P_j\in\R^{d\times r_j}$ has orthonormal principal directions as columns. The integer $r_j$ is the retained local rank. Define
\[
 s_j=\max\left\{n_j^{-1}\sum_{i=1}^{n_j}\|q_{ji}-\mu_j\|_2^2,10^{-6}\right\}.
\]
The locality weight, or gate, of model $j\in\mathcal J_b$ is
\begin{equation}
 g_j(z)=\frac{\exp\{-\|z-\mu_j\|_2^2/(2s_j)\}}
 {\sum_{\ell\in\mathcal J_b}\exp\{-\|z-\mu_\ell\|_2^2/(2s_\ell)\}}.
 \label{eq:gates}
\end{equation}
These gates are positive in exact arithmetic and sum to one within each partition. If $\Delta_q=\{a\in\R^q:a_i\geq0,\ \mathbf1_q^\top a=1\}$, their concatenation belongs to $\prod_{b=1}^{B_0}\Delta_{|\mathcal J_b|}$. There is no normalization across partitions. The weights describe proximity to reference regions; they are not assumed to be posterior probabilities of physical regimes.

The first two raw feature blocks are
\begin{equation}
 b_G(z)=\concat_j g_j(z),\qquad
 b_P(z)=\concat_j\{g_j(z)P_j^\top(z-\mu_j)\}.
 \label{eq:position}
\end{equation}
The first records membership and the second records local position. Their concatenation is denoted by $\psi_0(z)$. It can have more coordinates than the history itself, so no dimensionality-reduction claim is implied.

\subsection{Matched nonlinear reconstructions}
For $r_j>0$, let $S_j$ be the sample covariance of the projected references $P_j^\top q_{ji}$, computed with divisor $n_j-1$. Positive-variance principal directions make $S_j$ positive definite. The shared projected Gaussian kernel is
\begin{equation}
 \kappa_j(z,z')=\exp\left\{-\frac{[P_j^\top(z-z')]^\top S_j^{-1}[P_j^\top(z-z')]}{2r_j}\right\}.
 \label{eq:kernel}
\end{equation}
Let $[k_j(z)]_i=\kappa_j(z,q_{ji})$ and $[K_j]_{i\ell}=\kappa_j(q_{ji},q_{j\ell})$. The positive regularizer $\lambda_j$ is fitted from reference reconstruction on A, without forecast targets, using the iteration specified in the supplement. Set $C_j=(K_j+\lambda_j I_{n_j})^{-1}$. The KAHM, positive-kernel RBF, and affine reconstructions are
\begin{align}
 A_{K,j}(z)&=\frac{Q_j^\top C_j k_j(z)}{\mathbf1_{n_j}^\top C_j k_j(z)},&
 A_{R,j}(z)&=\frac{Q_j^\top k_j(z)}{\mathbf1_{n_j}^\top k_j(z)},\label{eq:reconstructions}\\
 A_{A,j}(z)&=\mu_j+P_jP_j^\top(z-\mu_j).\nonumber
\end{align}
The signed denominator in $A_{K,j}$ must be nonzero. The RBF denominator is positive for finite inputs. KAHM weights sum to one where defined, but can be negative; they are distinct from the nonnegative locality gates.

For $m\in\{K,R\}$, define the departure block
\begin{equation}
 d_m(z)=\concat_j\{g_j(z)[A_{m,j}(z)-A_{A,j}(z)]\}.
 \label{eq:departures}
\end{equation}
Each departure retains all $d$ ambient coordinates, including components outside the local principal subspace. Within a bank, K and R share the references, projected kernels, gates, affine projections, and departure width. Their raw difference is the inverse-Gram factor $C_j$ in the normalized reconstruction.

Numerical evaluation rescales each kernel vector by a common positive factor to avoid complete underflow. The KAHM reconstruction uses affine fallback when a denominator or finite-value check fails; in particular, it requires
\[
 \frac{|\mathbf1_{n_j}^\top C_jk_j(z)|}{\|C_jk_j(z)\|_1}\geq10^{-10}
\]
when the denominator of this test is nonzero. Affine fallback makes the corresponding raw departure zero. A rank-zero model also uses affine fallback, with $A_{A,j}(z)=\mu_j$. The supplement details the numerical safeguards and rank rules, and the source manifest identifies their implementation. The exact span results in Section~\ref{sec:theory} concern the mathematical maps where their denominators are nonzero, not a piecewise fallback implementation.

\subsection{Block transformations and affine readouts}
For each raw block $b$ of positive width $p$, let $a_b$ and $s_b$ be its A-fitted column means and population standard deviations, with every scale replaced by its maximum with $10^{-6}$. Division by $s_b$ is coordinatewise. Define
\begin{equation}
 \widetilde b(z)=T_p\bigl((b(z)-a_b)/s_b\bigr),\qquad
 T_p(v)=\frac{v}{\sqrt{1+\|v\|_2^2/p}}.
 \label{eq:transform}
\end{equation}
An empty block stays empty. Each block is transformed separately, and there is no subsequent transformation of their concatenation. The four readout inputs are
\begin{align}
 \chi_G&=\widetilde b_G,&
 \chi_P&=\concat(\widetilde b_G,\widetilde b_P),\nonumber\\
 \chi_K&=\concat(\widetilde b_G,\widetilde b_P,\widetilde d_K),&
 \chi_R&=\concat(\widetilde b_G,\widetilde b_P,\widetilde d_R).
 \label{eq:arms}
\end{align}
The membership block is identical across all four inputs, and the position block is identical across P, K, and R. Appending a block therefore leaves existing coordinates unchanged. The transformed memberships and signed enriched features are not generally simplex-valued.

For an arm $m$, a horizon $h$, and a fixed candidate bank, let $\chi_{m,i}$ and $y_{ih}$ be the A features and observed targets, and let $n_A$ be the number of retained A rows. The coefficients solve
\begin{equation}
 (\widehat a_{mh},\widehat w_{mh})=
 \arg\min_{a,w}\left\{\frac1{n_A}\sum_{i\in A}(y_{ih}-a-\chi_{m,i}^\top w)^2+\alpha\|w\|_2^2\right\},
 \quad \widehat f_{mh}(x)=\widehat a_{mh}+\chi_m(x)^\top\widehat w_{mh}.
 \label{eq:ridge}
\end{equation}
Here $\chi_m(x)$ includes the A-fitted standardization of $x$. The intercept is unpenalized. B selects $\alpha$ from $\{10^{-4},10^{-3},10^{-2},10^{-1},1\}$ separately for every arm, bank, and horizon. The target is the observed future response, not a residual or a successor feature vector. Although the coordinate sets are nested, fitted test loss need not decrease when coordinates are appended: the selected penalty and finite-sample estimation also change.
